%==========================================================
%  GLOSSARY OF KEY TERMS
%  ECON 204 — Contemporary Economic Issues: A Canadian Perspective
%==========================================================

\chapter*{Glossary of Key Terms}
\addcontentsline{toc}{chapter}{Glossary}
\markboth{Glossary}{Glossary}

\noindent This glossary defines key economic terms introduced in this textbook. Terms are listed alphabetically, with cross-references to the chapter in which each term first appears.

\vspace{12pt}

%----------------------------------------------------------
% A
%----------------------------------------------------------
\noindent\textbf{Adaptation (climate):} Adjustments in economic activity, infrastructure, and behaviour to reduce the harm caused by climate change that is already occurring or locked in. \textit{See} Chapter~\ref{ch:climate}.

\vspace{6pt}
\noindent\textbf{Aggregate Demand (AD):} The total quantity of goods and services demanded across the entire economy at each price level; the sum of consumption, investment, government expenditure, and net exports. \textit{See} Chapter~\ref{ch:inflation}.

\vspace{6pt}
\noindent\textbf{Aggregate Supply (AS):} The total quantity of goods and services that firms in an economy are willing and able to produce at each price level. The short-run aggregate supply (SRAS) slopes upward; the long-run aggregate supply (LRAS) is vertical at potential output. \textit{See} Chapter~\ref{ch:inflation}.

\vspace{6pt}
\noindent\textbf{Automation:} The replacement of human labour with machines, software, or robots in the performance of specific tasks. \textit{See} Chapter~\ref{ch:technology}.

%----------------------------------------------------------
% B
%----------------------------------------------------------
\vspace{6pt}
\noindent\textbf{Bank of Canada:} Canada's central bank, established in 1934 as a Crown corporation. It conducts monetary policy with the primary objective of maintaining low and stable inflation within a 1--3\% target range. \textit{See} Chapter~\ref{ch:inflation}.

\vspace{6pt}
\noindent\textbf{Built-in inflation:} Inflation arising from the expectation of future inflation; wage and price increases made in anticipation of continued price rises, creating a self-fulfilling wage-price spiral. \textit{See} Chapter~\ref{ch:inflation}.

%----------------------------------------------------------
% C
%----------------------------------------------------------
\vspace{6pt}
\noindent\textbf{Canada Carbon Rebate (CCR):} The payment made by the federal government to eligible Canadian households to return revenue from the consumer fuel charge. Paid quarterly as a flat amount per household. \textit{See} Chapter~\ref{ch:climate}.

\vspace{6pt}
\noindent\textbf{Canada Child Benefit (CCB):} A non-taxable, income-tested monthly benefit for families with children under 18. Launched in 2016 to replace previous child benefit programs; credited with significantly reducing child poverty in Canada. \textit{See} Chapter~\ref{ch:inequality}.

\vspace{6pt}
\noindent\textbf{Cap-and-trade:} An emissions reduction policy that sets a maximum (cap) on total greenhouse gas emissions and allows firms to buy and sell emissions permits. \textit{See} Chapter~\ref{ch:climate}.

\vspace{6pt}
\noindent\textbf{Capital gains:} The profit realized from selling an asset (real estate, shares, bonds) for more than its purchase price. \textit{See} Chapter~\ref{ch:inequality}.

\vspace{6pt}
\noindent\textbf{Carbon leakage:} The displacement of emissions-intensive production to jurisdictions with lower or no carbon prices as a result of domestic carbon pricing policy. \textit{See} Chapter~\ref{ch:climate}.

\vspace{6pt}
\noindent\textbf{Carbon tax:} A direct fee levied on greenhouse gas emissions, proportional to the carbon content of fossil fuels burned. A Pigouvian tax designed to internalize the externality cost of carbon emissions. \textit{See} Chapter~\ref{ch:climate}.

\vspace{6pt}
\noindent\textbf{Central bank independence:} The principle that a central bank has the authority to set monetary policy without direct political interference from the elected government. \textit{See} Chapter~\ref{ch:inflation}.

\vspace{6pt}
\noindent\textbf{Climate mitigation:} Actions that reduce greenhouse gas emissions or enhance carbon sinks, thereby slowing or reducing future climate change. \textit{See} Chapter~\ref{ch:climate}.

\vspace{6pt}
\noindent\textbf{Consumer Price Index (CPI):} A measure of the average change over time in the prices paid by urban consumers for a representative fixed basket of goods and services, produced monthly by Statistics Canada. \textit{See} Chapter~\ref{ch:inflation}.

\vspace{6pt}
\noindent\textbf{Core inflation:} Measures of inflation that exclude or reduce the influence of volatile food and energy prices, providing a cleaner signal of underlying inflationary trends. The Bank of Canada uses CPI-trim, CPI-median, and CPI-common. \textit{See} Chapter~\ref{ch:inflation}.

\vspace{6pt}
\noindent\textbf{Cost-push inflation:} Inflation caused by an increase in the cost of inputs to production (e.g., oil, wages), which reduces aggregate supply. \textit{See} Chapter~\ref{ch:inflation}.

%----------------------------------------------------------
% D
%----------------------------------------------------------
\vspace{6pt}
\noindent\textbf{Data network effects:} The process by which a platform with more users generates more data, improving its AI models and service quality, which attracts more users---a self-reinforcing competitive advantage. \textit{See} Chapter~\ref{ch:technology}.

\vspace{6pt}
\noindent\textbf{Deflation:} A sustained decrease in the general price level. Deflation can trap an economy in a downward spiral of postponed spending and rising real debt burdens. \textit{See} Chapter~\ref{ch:inflation}.

\vspace{6pt}
\noindent\textbf{Demand-pull inflation:} Inflation caused by an increase in aggregate demand, resulting in ``too much money chasing too few goods.'' \textit{See} Chapter~\ref{ch:inflation}.

\vspace{6pt}
\noindent\textbf{Development charges:} Fees levied by municipal governments on builders and developers to fund the infrastructure (roads, water, parks, transit) required to support new residential and commercial development. \textit{See} Chapter~\ref{ch:housing}.

%----------------------------------------------------------
% E
%----------------------------------------------------------
\vspace{6pt}
\noindent\textbf{Exclusionary zoning:} Land-use regulations that prohibit multi-unit residential buildings (apartments, duplexes, townhouses) in residential zones, permitting only detached single-family homes. A primary driver of housing supply constraints in Canadian cities. \textit{See} Chapter~\ref{ch:housing}.

%----------------------------------------------------------
% F
%----------------------------------------------------------
\vspace{6pt}
\noindent\textbf{Fiscal policy:} The use of government spending and taxation to influence aggregate demand and economic output. Controlled by the elected government through the Department of Finance Canada. \textit{See} Chapter~\ref{ch:introduction}.

\vspace{6pt}
\noindent\textbf{Fisher equation:} The relationship between nominal interest rates, real interest rates, and inflation: $r \approx i - \pi$, where $r$ is the real rate, $i$ is the nominal rate, and $\pi$ is inflation. \textit{See} Chapter~\ref{ch:inflation}.

\vspace{6pt}
\noindent\textbf{Free Prior and Informed Consent (FPIC):} A principle under the United Nations Declaration on the Rights of Indigenous Peoples (UNDRIP) stating that Indigenous peoples have the right to give or withhold consent to resource development on their traditional territories. \textit{See} Chapter~\ref{ch:inequality}.

\vspace{6pt}
\noindent\textbf{Frictional unemployment:} Short-term unemployment arising from the process of matching workers with jobs; inherent in any dynamic labour market. \textit{See} Chapter~\ref{ch:labour}.

%----------------------------------------------------------
% G
%----------------------------------------------------------
\vspace{6pt}
\noindent\textbf{Gender wage gap:} The difference between median earnings of men and women in the labour market, typically expressed as women's earnings as a percentage of men's. \textit{See} Chapter~\ref{ch:labour}.

\vspace{6pt}
\noindent\textbf{General-purpose technology (GPT):} An innovation that is pervasive across many sectors, capable of continuous improvement, and able to generate complementary innovations. Examples include steam power, electricity, computers, and AI. \textit{See} Chapter~\ref{ch:technology}.

\vspace{6pt}
\noindent\textbf{Gig economy:} A labour market characterized by short-term contracts and freelance work, often mediated by digital platforms. Workers are typically classified as independent contractors. \textit{See} Chapter~\ref{ch:labour}.

\vspace{6pt}
\noindent\textbf{Gini coefficient:} A summary measure of income (or wealth) inequality derived from the Lorenz curve. Ranges from 0 (perfect equality) to 1 (perfect inequality). \textit{See} Chapter~\ref{ch:inequality}.

\vspace{6pt}
\noindent\textbf{Green industrial policy:} Government interventions (subsidies, tax credits, public investment, mandates) designed to accelerate the deployment of clean energy technologies and the transition to a low-carbon economy. \textit{See} Chapter~\ref{ch:climate}.

\vspace{6pt}
\noindent\textbf{Gross Domestic Product (GDP):} The total market value of all final goods and services produced within a country's borders in a given period (typically one year). \textit{See} Chapter~\ref{ch:introduction}.

%----------------------------------------------------------
% H
%----------------------------------------------------------
\vspace{6pt}
\noindent\textbf{Housing affordability:} The relationship between housing costs (prices, rents) and household income. Commonly measured by the price-to-income ratio or the share of income spent on shelter. \textit{See} Chapter~\ref{ch:housing}.

%----------------------------------------------------------
% I
%----------------------------------------------------------
\vspace{6pt}
\noindent\textbf{Imputed rent:} The estimated rental value of owner-occupied housing---the income an owner would receive by renting their property, or equivalently, the rent they save by owning. \textit{See} Chapter~\ref{ch:housing}.

\vspace{6pt}
\noindent\textbf{Income quintile:} One-fifth of the population ranked by income. \textit{See} Chapter~\ref{ch:inequality}.

\vspace{6pt}
\noindent\textbf{Inflation:} A sustained increase in the general price level of goods and services in an economy over time. \textit{See} Chapter~\ref{ch:inflation}.

\vspace{6pt}
\noindent\textbf{Inflation targeting:} A monetary policy framework in which the central bank publicly commits to maintaining inflation at a specific numerical target (in Canada, 2\% within a 1--3\% range). \textit{See} Chapter~\ref{ch:inflation}.

%----------------------------------------------------------
% K
%----------------------------------------------------------
\vspace{6pt}
\noindent\textbf{Kuznets hypothesis:} Simon Kuznets's 1955 prediction that income inequality follows an inverted-U pattern as an economy develops: rising during early industrialization and falling as development matures. Empirically challenged by the experience of advanced economies since 1980. \textit{See} Chapter~\ref{ch:inequality}.

%----------------------------------------------------------
% L
%----------------------------------------------------------
\vspace{6pt}
\noindent\textbf{Labour force participation rate:} The percentage of the working-age population (15+) that is either employed or actively seeking work. \textit{See} Chapter~\ref{ch:labour}.

\vspace{6pt}
\noindent\textbf{Labour productivity:} Output per unit of labour input; typically measured as real GDP per hour worked. \textit{See} Chapter~\ref{ch:labour}.

\vspace{6pt}
\noindent\textbf{Lorenz curve:} A graphical representation of the income (or wealth) distribution, plotting the cumulative share of total income against the cumulative share of the population ranked from poorest to richest. \textit{See} Chapter~\ref{ch:inequality}.

\vspace{6pt}
\noindent\textbf{Low Income Cut-Off (LICO):} Statistics Canada's original poverty measure, based on the share of income devoted to food, shelter, and clothing. No longer promoted as an official measure but still published. \textit{See} Chapter~\ref{ch:inequality}.

\vspace{6pt}
\noindent\textbf{Low Income Measure (LIM):} A relative poverty measure defined as income below 50\% of median adjusted household income; the standard OECD cross-country poverty measure. \textit{See} Chapter~\ref{ch:inequality}.

%----------------------------------------------------------
% M
%----------------------------------------------------------
\vspace{6pt}
\noindent\textbf{Marginal Revenue Product of Labour ($MRP_L$):} The additional revenue generated by hiring one more unit of labour; the basis for the firm's labour demand curve. \textit{See} Chapter~\ref{ch:labour}.

\vspace{6pt}
\noindent\textbf{Market Basket Measure (MBM):} Canada's official poverty line since 2019. Measures the cost of a specific basket of goods and services needed for a modest but adequate standard of living, with regional variation. \textit{See} Chapter~\ref{ch:inequality}.

\vspace{6pt}
\noindent\textbf{Minimum wage:} A government-mandated floor on wages below which employers cannot legally pay workers. Set provincially in Canada, with a separate federal rate for federally regulated industries. \textit{See} Chapter~\ref{ch:labour}.

\vspace{6pt}
\noindent\textbf{Monetary policy:} Actions taken by a central bank to influence interest rates, the money supply, and ultimately inflation and economic activity. In Canada, conducted by the Bank of Canada through the overnight rate. \textit{See} Chapter~\ref{ch:introduction}.

\vspace{6pt}
\noindent\textbf{Monetary policy transmission:} The process by which a change in the central bank's policy rate affects economic variables such as mortgage rates, the exchange rate, credit conditions, asset prices, and ultimately inflation. \textit{See} Chapter~\ref{ch:inflation}.

\vspace{6pt}
\noindent\textbf{Monopsony:} A labour market in which there is only one buyer of labour (one employer), giving that employer market power to set wages below the competitive level. \textit{See} Chapter~\ref{ch:labour}.

\vspace{6pt}
\noindent\textbf{Mortgage stress test:} OSFI's regulatory requirement (B-20) that mortgage borrowers must qualify at the higher of their contract rate plus 2 percentage points, or 5.25\%, ensuring borrowers have a buffer against rate increases. \textit{See} Chapter~\ref{ch:inflation}.

%----------------------------------------------------------
% N
%----------------------------------------------------------
\vspace{6pt}
\noindent\textbf{Natural rate of unemployment:} The rate of unemployment consistent with full employment---the sum of frictional and structural unemployment; the rate toward which the economy gravitates in the long run. \textit{See} Chapter~\ref{ch:labour}.

\vspace{6pt}
\noindent\textbf{Negative externality:} A cost imposed on a third party (not the buyer or seller) by an economic transaction that is not reflected in the market price. \textit{See} Chapter~\ref{ch:climate}.

\vspace{6pt}
\noindent\textbf{Network effects:} The phenomenon whereby the value of a product or service to any one user increases as more users adopt it. Network effects create tendencies toward winner-take-all market outcomes. \textit{See} Chapter~\ref{ch:technology}.

\vspace{6pt}
\noindent\textbf{Normative economics:} Economic analysis that makes value judgements about what \textit{should} happen, rather than what \textit{is} or \textit{will} happen. Contrasts with positive economics. \textit{See} Chapter~\ref{ch:introduction}.

%----------------------------------------------------------
% O
%----------------------------------------------------------
\vspace{6pt}
\noindent\textbf{Opportunity cost:} The value of the next-best alternative foregone when a choice is made. The true economic cost of any decision. \textit{See} Chapter~\ref{ch:introduction}.

\vspace{6pt}
\noindent\textbf{Overnight rate:} The Bank of Canada's policy interest rate---the target for the rate at which major financial institutions lend to each other overnight. Changes in the overnight rate transmit through the financial system to affect mortgage rates, business loans, and the exchange rate. \textit{See} Chapter~\ref{ch:inflation}.

%----------------------------------------------------------
% P
%----------------------------------------------------------
\vspace{6pt}
\noindent\textbf{Pigouvian tax:} A tax set equal to the marginal external cost of a negative externality, designed to align private incentives with social costs. Named after economist Arthur C. Pigou. The carbon tax is a Pigouvian tax. \textit{See} Chapter~\ref{ch:climate}.

\vspace{6pt}
\noindent\textbf{Platform economy:} A market structure in which a digital intermediary (platform) connects two or more groups of users (e.g., buyers and sellers, drivers and riders) and generates value primarily from network effects and data. \textit{See} Chapter~\ref{ch:technology}.

\vspace{6pt}
\noindent\textbf{Positive economics:} Economic analysis that describes and explains economic phenomena without making value judgements; empirical and testable claims about how the economy works. \textit{See} Chapter~\ref{ch:introduction}.

\vspace{6pt}
\noindent\textbf{Poverty:} A condition in which a household or individual lacks sufficient income or resources to meet a defined minimum standard of living. May be measured as absolute poverty (inability to meet basic needs) or relative poverty (income below a fraction of the median). \textit{See} Chapter~\ref{ch:inequality}.

\vspace{6pt}
\noindent\textbf{Price elasticity of supply:} The percentage change in quantity supplied resulting from a 1\% increase in price. Inelastic supply ($E_s < 1$) means supply responds weakly to price. \textit{See} Chapter~\ref{ch:housing}.

\vspace{6pt}
\noindent\textbf{Productivity paradox:} The observation that investment in information and communication technologies does not immediately produce measurable aggregate productivity gains---the lag between GPT adoption and full productivity impact. \textit{See} Chapter~\ref{ch:technology}.

%----------------------------------------------------------
% R
%----------------------------------------------------------
\vspace{6pt}
\noindent\textbf{Real interest rate:} The nominal interest rate adjusted for inflation: $r \approx i - \pi$. The real interest rate measures the true cost of borrowing in terms of purchasing power. \textit{See} Chapter~\ref{ch:inflation}.

\vspace{6pt}
\noindent\textbf{Rent control:} Government regulation limiting the rent a landlord may charge or the permissible annual rent increase for a residential unit. \textit{See} Chapter~\ref{ch:housing}.

\vspace{6pt}
\noindent\textbf{Routine vs.\ non-routine tasks:} In the task-based model of automation, routine tasks are governed by explicit, codifiable rules (most automatable); non-routine tasks require judgment, creativity, or physical dexterity in variable environments (less automatable). \textit{See} Chapter~\ref{ch:technology}.

%----------------------------------------------------------
% S
%----------------------------------------------------------
\vspace{6pt}
\noindent\textbf{Skill-biased technological change (SBTC):} Technological change that increases relative demand for skilled (typically higher-educated) workers and decreases relative demand for unskilled workers, widening the wage premium for education. \textit{See} Chapter~\ref{ch:inequality}.

\vspace{6pt}
\noindent\textbf{Social cost of carbon (SCC):} The economic value of the damage caused by emitting one additional tonne of CO$_2$-equivalent into the atmosphere, including all present and future climate harms. \textit{See} Chapter~\ref{ch:climate}.

\vspace{6pt}
\noindent\textbf{Stagflation:} The simultaneous occurrence of high inflation and low economic output (or recession). A particularly challenging macroeconomic condition to address with conventional policy tools. \textit{See} Chapter~\ref{ch:inflation}.

\vspace{6pt}
\noindent\textbf{Structural unemployment:} Long-term unemployment resulting from a mismatch between workers' skills and available jobs, or from geographic mismatch. \textit{See} Chapter~\ref{ch:labour}.

%----------------------------------------------------------
% T
%----------------------------------------------------------
\vspace{6pt}
\noindent\textbf{Task-based model of automation:} An economic framework that classifies tasks (rather than whole jobs) by their susceptibility to automation, distinguishing routine from non-routine and cognitive from manual tasks. \textit{See} Chapter~\ref{ch:technology}.

\vspace{6pt}
\noindent\textbf{Time inconsistency problem:} The tendency for a policy-maker who has committed to a future policy to have incentives to deviate from that commitment once the relevant time arrives. The problem underlies the case for central bank independence. \textit{See} Chapter~\ref{ch:inflation}.

\vspace{6pt}
\noindent\textbf{Trade-off:} The situation in which choosing more of one thing means getting less of another. A fundamental concept in economics reflecting scarcity. \textit{See} Chapter~\ref{ch:introduction}.

%----------------------------------------------------------
% U
%----------------------------------------------------------
\vspace{6pt}
\noindent\textbf{Unemployment rate:} The percentage of the labour force that is without work but actively seeking employment. \textit{See} Chapter~\ref{ch:labour}.

\vspace{6pt}
\noindent\textbf{Union density:} The share of paid workers who are members of a trade union. \textit{See} Chapter~\ref{ch:labour}.

\vspace{6pt}
\noindent\textbf{UNDRIP (United Nations Declaration on the Rights of Indigenous Peoples):} An international human rights instrument adopted by the UN General Assembly in 2007. Canada passed legislation in 2021 (Bill C-15) to align Canadian laws with UNDRIP. \textit{See} Chapter~\ref{ch:inequality}.

%----------------------------------------------------------
% W
%----------------------------------------------------------
\vspace{6pt}
\noindent\textbf{Wage polarization:} The pattern of employment and wage growth observed in advanced economies: expansion at the top (high-skill, high-wage) and bottom (low-skill, low-wage) of the distribution, with contraction in the middle (middle-skill, routine occupations). Attributed to skill-biased technological change and automation. \textit{See} Chapter~\ref{ch:technology}.

\vspace{6pt}
\noindent\textbf{Wealth tax:} An annual levy on the total net worth (assets minus liabilities) of individuals above a specified threshold, as distinct from taxes on income flows. \textit{See} Chapter~\ref{ch:inequality}.

\vspace{6pt}
\noindent\textbf{Winner-take-all market:} A market in which one or a few firms capture most of the market share due to strong network effects, economies of scale, or first-mover advantages. Common in digital platform markets. \textit{See} Chapter~\ref{ch:technology}.
